% Section 4 -- Method (P1)
% Every setting here comes from src/qaoa.py, src/problem.py or experiments/E1_depth.py.

\section{Method}
\label{sec:method}

This section describes how we build, optimise and evaluate the QAOA circuit for the cost Hamiltonian $H_C$ of Eq.~\eqref{eq:hamiltonian}.
Section~\ref{sec:qaoa-background} introduced the algorithm; here we fix the choices that the results depend on.

\subsection{Circuit}
\label{sec:circuit}

The circuit uses one qubit per job, with qubit $i$ holding job $i$.
It starts with a Hadamard gate on every qubit to prepare $\lvert + \rangle^{\otimes n}$ and then applies $p$ layers.
In layer $\ell$ the cost unitary $e^{-i\gamma_\ell H_C}$ factorises exactly into one two-qubit term per coupling, because all terms of $H_C$ commute.
Each term is compiled as
\begin{equation}
  e^{-i \gamma J_{ij} Z_i Z_j} = \mathrm{CX}_{ij}\; R_Z^{(j)}(2\gamma J_{ij})\; \mathrm{CX}_{ij},
  \label{eq:zz-gate}
\end{equation}
and the mixer $e^{-i\beta_\ell \sum_i X_i}$ is one $R_X(2\beta_\ell)$ rotation on every qubit.
The constant term of $H_C$ is a global phase and is not applied.
A layer therefore contains $n(n-1)$ CX gates, $n(n-1)/2$ $R_Z$ gates and $n$ $R_X$ gates, so the gate count grows as $n^2$ per layer and linearly in $p$.

The $2p$ angles are stored as one vector $\boldsymbol\theta = (\gamma_1, \dots, \gamma_p, \beta_1, \dots, \beta_p)$.
Because $H_C$ is real, the angles $\boldsymbol\theta$ and $-\boldsymbol\theta$ give complex-conjugate states with identical measurement probabilities.
Before we compare or average angles across instances, we therefore flip the sign of any vector with $\gamma_1 < 0$, so that all of them lie on the $\gamma_1 \geq 0$ branch.

\subsection{Objective and optimiser}
\label{sec:optimiser}

The optimiser minimises the exact expectation $E(\boldsymbol\theta) = \langle \boldsymbol\theta \rvert H_C \lvert \boldsymbol\theta \rangle$, which the simulator computes directly from the statevector.
Using the exact value instead of a shot estimate removes sampling noise from the optimisation, so the results isolate the behaviour of QAOA itself; on hardware, $E$ would have to be estimated from measurements.
Because $H_C$ is scaled so that $\max |J_{ij}| = 1$ (Section~\ref{sec:normalisation}), the same angle scale applies to every instance.

We use COBYLA, a gradient-free trust-region method, as implemented in SciPy, with an initial step of $0.2$\,rad and a budget of 500 energy evaluations per optimisation.
Gradient-free methods are the usual choice for QAOA because each gradient would cost $4p$ additional circuit evaluations.
We record the number of evaluations each run uses, since this is the main driver of simulation time.

\subsection{Initial angles}
\label{sec:init}

The energy landscape has many local minima, so the starting point matters.
We use four strategies; E1 uses the first, and E2 compares all four.

\begin{description}
  \item[Ramp ($p = 1$ only in E1).]
  A linear ramp that imitates a discretised annealing schedule: $\gamma_\ell = 0.75\,t_\ell$ and $\beta_\ell = 0.75\,(1 - t_\ell)$ with $t_\ell = (\ell - \tfrac12)/p$.
  \item[INTERP.]
  The method of Zhou et al.~\cite{zhou2020}.
  The optimised angles at depth $p$ are linearly interpolated onto $p + 1$ points and used as the starting point at depth $p + 1$:
  \begin{equation}
    \gamma^{(p+1)}_{\ell} = \frac{\ell - 1}{p}\, \gamma^{(p)}_{\ell - 1} + \frac{p - \ell + 1}{p}\, \gamma^{(p)}_{\ell},
    \qquad \ell = 1, \dots, p + 1,
  \end{equation}
  with $\gamma^{(p)}_0 = \gamma^{(p)}_{p+1} = 0$, and the same rule for $\boldsymbol\beta$.
  Each instance is therefore optimised as a chain $p = 1, 2, \dots, 6$, and depth $p + 1$ starts from a state at least as good as the optimum at depth $p$.
  \item[Random.]
  $\gamma_\ell$ uniform in $[0, \pi]$ and $\beta_\ell$ uniform in $[0, \pi/2]$; we keep the best of five starts.
  \item[Transfer.]
  The element-wise median of the optimised (sign-canonicalised) angles over the ten $n = 8$ instances of E1, used either as a starting point or directly with no optimisation.
  Optimal QAOA angles are known to concentrate across instances of the same problem class, so angles found on small instances are often good on larger ones.
\end{description}

\subsection{Sampling and metrics}
\label{sec:metrics}

After optimisation the circuit is measured 10{,}000 times, and each bitstring is mapped to a schedule with $z_i = 1 - 2x_i$.
Every instance used in E1 and E2 has at most 20 jobs, so we compute its exact minimum $C_{\min}$, maximum $C_{\max}$ and the full set of optimal bitstrings by brute force over all $2^n$ schedules.
For $n > 24$ (E5) the minimum comes from CP-SAT and $C_{\max}$ from a closed form.
We report four metrics.

\begin{enumerate}
  \item \textbf{Approximation ratio} of the mean cost,
  \begin{equation}
    r = \frac{C_{\max} - \langle C \rangle}{C_{\max} - C_{\min}},
    \label{eq:approx-ratio}
  \end{equation}
  with $\langle C \rangle = s\,E(\boldsymbol\theta)$.
  It is 1 for a state concentrated on optimal schedules and 0 for the worst schedule.
  Because the cost is quadratic and its range is set by the worst schedule, the uniform superposition (random guessing) already scores $r \approx 0.71$--$0.73$ on our instances; this is the $p = 0$ reference.
  \item \textbf{Success probability} $P(\text{opt})$: the fraction of shots that return \emph{any} optimal bitstring (Section~\ref{sec:symmetry}).
  \item \textbf{Best sample}: the lowest-cost schedule among the 10{,}000 shots, reported as its own approximation ratio. This is the answer QAOA would return in practice, and it is the like-for-like comparison with classical solvers, which also return one schedule.
  \item \textbf{Violations and makespan} of the best sample, so that a schedule that trades a conflict for balance is always visible.
\end{enumerate}

Every run also records its wall time and number of energy evaluations.
Runs are identified by (experiment, instance, depth, initialisation, backend, precision, noise level, seed) and appended to a CSV file as soon as they finish, so an interrupted sweep resumes where it stopped and every figure is regenerated from these files by a single script.

\subsection{Validation}
\label{sec:validation}

Before running any experiment we checked the implementation against exact results with an automated test suite (40 tests).
With all angles at zero the circuit returns the uniform state and $\langle C \rangle = \sum_i w_i^2$; the energy from the simulator matches the probability-weighted cost computed from the exact statevector; the sampled mean cost matches the exact expectation within shot noise; every bitstring and its mirror image have the same probability (Section~\ref{sec:symmetry}); and the qubit ordering maps bit $i$ to job $i$.
As end-to-end checks, $p = 3$ on a 6-job instance reaches $r > 0.9$ and samples the optimum more than five times as often as random guessing, and optimised QAOA beats random sampling on both metrics at $n = 10$.
For the noise study (E6), every noiseless density-matrix result is checked against the statevector energy and the run stops if they disagree.
